% =========================================================
% NeurIPS 2026 E&D Blueprint
% Replace "neurips_2026" below with the exact style filename
% shipped in the official NeurIPS 2026 template zip if needed.
% =========================================================
\documentclass{article}

 % for drafting / arXiv
% \usepackage{neurips_2026}         % for anonymous submission
% \usepackage[final]{neurips_2026}  % for camera-ready

\usepackage[utf8]{inputenc}
\usepackage[T1]{fontenc}
\usepackage{microtype}
\usepackage{url}
\usepackage{booktabs}
\usepackage{amsmath,amssymb,amsfonts,amsthm}
\usepackage{mathtools}
\usepackage{nicefrac}
\usepackage{xcolor}
\usepackage{graphicx}
\usepackage{subcaption}
\usepackage{multirow}
\usepackage{array}
\usepackage{enumitem}
\usepackage{siunitx}
\usepackage{hyperref}
\usepackage[nameinlink,capitalize]{cleveref}

% -------------------------------
% Macros
% -------------------------------
\newcommand{\Recall}{\textsc{Recall}}
\newcommand{\ADP}{\textsc{ADP}}
\newcommand{\OH}{\textsc{OpenHands}}
\newcommand{\etal}{\textit{et al.}}
\newcommand{\todo}[1]{\textcolor{red}{[TODO: #1]}}
\newcommand{\note}[1]{\textcolor{blue}{[NOTE: #1]}}

\DeclareMathOperator*{\argmax}{arg\,max}

% -------------------------------
% Title options
% -------------------------------
% Option A: evaluation-first
\title{Does Structured Trajectory Memory Improve Agent Retrieval?\\
Evaluating \ADP{} as a Representation for Experience-Augmented Agents}

% Option B: more direct
% \title{From Training Format to Memory Substrate: Evaluating \ADP{} for Trajectory Retrieval in Agentic RAG}

% Option C: stronger emphasis on retrieval-control
% \title{When Should Agents Retrieve Past Experience? Evaluating \ADP{}-Based Procedural Memory for Agentic RAG}

\author{Anonymous Authors\\
Paper under double-blind review}

\begin{document}

\maketitle

% =========================================================
% Abstract
% =========================================================
\begin{abstract}
Agent Data Protocol (\ADP{}) was originally introduced as a standardized and expressive schema for
representing heterogeneous agent trajectories, with demonstrated value for offline agent training and
fine-tuning. In this paper, we study a different question: whether \ADP{} is also a better
\emph{retrieval representation} for experience-augmented agents.
We consider a retrieval setting in which an agent can query a memory store of prior trajectories during
live execution, retrieve candidate past experiences, and optionally inject them back into context.
Our central hypothesis is that \ADP{} preserves procedural structure---tool usage, action types,
environment observations, and interaction flow---that is lost in weaker memory representations such as
flattened text logs or short natural-language summaries.
To test this, we build a reproducible evaluation pipeline that converts agent sessions into \ADP{},
indexes them as memory items, and compares \ADP{}-based retrieval against unstructured baselines across
tool-using agent tasks.
We evaluate not only task success, but also retrieval usefulness, injection harm, search cost, and
robustness under noisy or stale memory.
Our results aim to answer three questions:
(1) whether \ADP{} improves procedural memory retrieval,
(2) which structural components of \ADP{} matter most, and
(3) when retrieval is beneficial versus harmful in live agent execution.
Taken together, the paper positions \ADP{} not only as a training interlingua, but also as a candidate
memory substrate for retrieval-augmented agents.
\end{abstract}

% =========================================================
% 1. Introduction
% =========================================================
\section{Introduction}
Large language model agents increasingly operate through multi-step interactions with tools,
environments, and users.
In such settings, purely parametric memory is often insufficient:
agents may repeatedly encounter similar subproblems, failure modes, and tool-use patterns across tasks.
A natural idea is to augment a live agent with a memory store of prior trajectories and retrieve relevant
experience when needed.

However, the effectiveness of such retrieval depends not only on the retrieval algorithm but also on the
\emph{representation} of the stored experience itself.
Most existing memory implementations flatten trajectories into plain text, short summaries, or ad hoc log
formats.
These formats are convenient, but they often discard procedural structure:
tool invocations, typed actions, environment feedback, and intermediate failure signals may be lost or
blurred into unstructured text.
This makes it unclear whether current retrieval-based agent memory is underperforming because of weak
retrieval policies, weak storage representations, or both.

This paper investigates whether \ADP{}, originally proposed as a standardized representation for offline
agent training data, is also a better representation for retrieval-based agent memory.
The key intuition is simple:
if memory items preserve the procedural structure of past experience, then retrieval and post-retrieval
decision-making should be able to exploit more than superficial semantic overlap.

We study this question in a controlled setting built around a memory system called \Recall{}.
In \Recall{}, a live agent can query a trajectory store, retrieve candidate prior sessions, and decide
whether to inject nothing, a compact summary, or a fuller memory artifact into context.
A key design constraint is that some signals needed for relevance estimation, especially semantic
similarity to stored trajectories, are only available after retrieval has been executed.
This motivates a hierarchical retrieval-control view in which local search decisions and post-search
injection decisions are separated.

The paper makes the following contributions:
\begin{enumerate}[leftmargin=*,itemsep=2pt,topsep=2pt]
    \item We formulate \emph{trajectory representation for agent memory} as a scientific evaluation problem,
    rather than a purely systems problem.
    \item We build a reproducible evaluation pipeline that converts agent sessions into \ADP{} memory items,
    stores them in a searchable memory substrate, and supports controlled comparison against weaker
    representations.
    \item We compare \ADP{} against unstructured baselines such as flattened logs, message-only traces, and
    natural-language summaries.
    \item We analyze not only whether retrieval helps, but also when it hurts, which parts of \ADP{} matter,
    and how retrieval-control interacts with memory quality.
\end{enumerate}

Our goal is not to claim that \ADP{} universally solves agent memory.
Rather, the goal is to determine whether structured, typed trajectories provide a measurable advantage as
memory units for retrieval-augmented agent execution.

% =========================================================
% 2. Research Questions and Hypotheses
% =========================================================
\section{Research Questions and Hypotheses}
We center the paper around the following research questions.

\paragraph{RQ1.}
Does \ADP{}-based memory retrieval improve downstream agent performance compared with weaker trajectory
representations such as flattened text logs, message-only traces, or summary memory?

\paragraph{RQ2.}
Which components of \ADP{} matter most for retrieval effectiveness:
typed actions, tool arguments, observations, metadata, or their interaction?

\paragraph{RQ3.}
Does \ADP{} help more on procedural and tool-heavy tasks than on primarily conversational tasks?

\paragraph{RQ4.}
When is retrieval beneficial, and when is it harmful, in the presence of noisy, stale, or weakly matched
memory?

We test the following hypotheses.

\paragraph{H1.}
\ADP{} outperforms unstructured memory representations on tasks where procedural reuse matters.

\paragraph{H2.}
The gain from \ADP{} comes primarily from preserving typed tool/action structure and environment
observations, not merely from storing more tokens.

\paragraph{H3.}
Naive memory usage is often harmful or wasteful; the benefit of structured memory is largest when
combined with controlled retrieval and selective injection.

% =========================================================
% 3. Background and Related Work
% =========================================================
\section{Background and Related Work}

\subsection{\ADP{} as a trajectory standard}
Briefly summarize the original \ADP{} paper here:
\begin{itemize}[leftmargin=*,itemsep=2pt,topsep=2pt]
    \item motivation: heterogeneous agent datasets are difficult to combine;
    \item design principles: simplicity, standardization, expressiveness;
    \item schema: trajectory id, content sequence, flexible details metadata;
    \item action types: API actions, code actions, message actions;
    \item observation types: text observations, web observations;
    \item original demonstrated use case: offline fine-tuning / training.
\end{itemize}
\textbf{Important framing sentence:}
The original \ADP{} work studies \emph{training-time} utility; our paper studies \emph{retrieval-time}
utility.

\subsection{Memory and retrieval in LLM agents}
Review agent memory systems, trajectory retrieval, episodic memory, procedural memory,
experience replay for agents, and agentic RAG.
Emphasize the gap:
most prior work studies retrieval policy or memory storage, but not the effect of the
\emph{trajectory representation itself}.

\subsection{Evaluation as a scientific object}
Because this paper targets the E\&D track, make the evaluation contribution explicit.
Argue that representation choice can change what retrieval claims are even measurable,
and therefore belongs in the science of evaluation.

\subsection{Positioning of this paper}
State clearly:
This paper is \emph{not} proposing \ADP{} itself.
It asks whether \ADP{}, originally justified for training, also serves as a superior memory substrate for
retrieval-based agents.

% =========================================================
% 4. Problem Setup
% =========================================================
\section{Problem Setup}

We consider a live agent interacting with a user and an external environment over discrete decision steps
$t = 1,2,\ldots$.
At each step, the agent may emit a message, call a tool, receive an observation, or replan.
Alongside this live execution, the system maintains a memory store $\mathcal{M}$ of prior trajectories.

A memory item $m \in \mathcal{M}$ is derived from a past session and stored under one of several
representations:
\[
r \in \mathcal{R}
=
\{
\texttt{flat-text},
\texttt{message-only},
\texttt{summary},
\texttt{ADP},
\texttt{ADP-lite}
\}.
\]

At selected steps, the live session may issue a retrieval query $q_t$, obtain a candidate set
\[
C_t = \{c_1,\dots,c_k\},
\]
and optionally inject a memory-derived artifact into the agent context.

We separate the problem into two layers:
\begin{enumerate}[leftmargin=*,itemsep=2pt,topsep=2pt]
    \item \textbf{Representation quality:}
    does the chosen stored format make relevant experience easier to retrieve and use?
    \item \textbf{Retrieval-control quality:}
    given a representation, when should the system search, and when should it inject?
\end{enumerate}

This paper focuses primarily on the first question and treats retrieval-control as either
(a) a controlled baseline factor, or
(b) a secondary analysis dimension.

\subsection{Evaluation targets}
We evaluate memory systems along five axes:
\begin{enumerate}[leftmargin=*,itemsep=2pt,topsep=2pt]
    \item task success,
    \item retrieval usefulness,
    \item injection harm,
    \item search/context cost,
    \item robustness to memory noise and staleness.
\end{enumerate}

\subsection{Why representation matters}
Add a short intuitive example here:
two past sessions may share procedural structure while differing lexically.
A plain-text or summary representation may miss this similarity, while \ADP{} may preserve it through
typed tool calls, action ordering, and environment feedback.

% =========================================================
% 5. System Overview: Recall
% =========================================================
\section{System Overview: \Recall{}}

\subsection{Current implementation}
Describe only what is actually implemented.
This section should be factual and conservative.

\paragraph{Session normalization.}
We convert \OH{} session logs into \ADP{}-shaped trajectories.
The converter maps user messages, agent messages, terminal actions, file editor actions,
tool-related events, and environment observations into typed \ADP{} items.

\paragraph{Embedding and storage.}
Each converted trajectory is embedded and stored as a single indexed memory item together with metadata
and the full \ADP{} content.

\paragraph{Memory backend.}
The current prototype uses a vector store for approximate similarity search over prior trajectories.

\paragraph{Important scope boundary.}
The current repository implements offline trajectory conversion and ingestion.
Any online retrieval/injection loop used in experiments must be described separately and clearly marked as
part of the experimental harness.

\subsection{Memory representations compared}
Define each representation formally.

\paragraph{Flat-text memory.}
Flatten the entire trajectory into a single text string.

\paragraph{Message-only memory.}
Retain only user and agent natural-language messages; discard tool/action structure.

\paragraph{Summary memory.}
Compress the trajectory into a short natural-language summary of fixed budget.

\paragraph{\ADP{} memory.}
Store the typed action-observation trajectory and metadata as structured memory.

\paragraph{\ADP{}-lite memory.}
Ablated structured variant that removes selected components, e.g.\ tool kwargs or observations.

\subsection{Optional retrieval-control layer}
If you include the trigger, keep it secondary.
Do \emph{not} make the trigger the main contribution unless experiments justify it.

A minimal paragraph could say:
\begin{quote}
For part of our analysis, we consider a hierarchical retrieval-control policy in which the live wrapper
first decides whether to search, after which post-retrieval signals determine whether to inject no
memory, a compact summary, or a fuller artifact.
This decomposition reflects the fact that some relevance signals are only available after candidate
retrieval.
\end{quote}

% =========================================================
% 6. Experimental Design
% =========================================================
\section{Experimental Design}

\subsection{Task families}
Describe the task suite.
A strong structure is:
\begin{itemize}[leftmargin=*,itemsep=2pt,topsep=2pt]
    \item \textbf{Procedural/tool-heavy tasks:} coding, debugging, file editing, web workflows.
    \item \textbf{Mixed tasks:} tasks with both language interaction and tool use.
    \item \textbf{Low-procedural/conversational controls:} tasks where memory structure should matter less.
\end{itemize}

\subsection{Memory corpus construction}
Explain where stored trajectories come from.
Possible sources:
\begin{itemize}[leftmargin=*,itemsep=2pt,topsep=2pt]
    \item historical \OH{} sessions,
    \item replayed benchmark trajectories,
    \item held-out prior episodes from the same domain,
    \item synthetic perturbations for noise/staleness analysis.
\end{itemize}

\paragraph{Train/validation/test split for memory.}
State how memory items and evaluation tasks are separated.
This is crucial.
Avoid leakage between the memory bank and target episodes.

\subsection{Compared methods}
Your core table should compare:

\begin{enumerate}[label=\textbf{M\arabic*:},leftmargin=*,itemsep=2pt,topsep=2pt]
    \item No memory.
    \item Flat-text memory.
    \item Message-only memory.
    \item Summary memory.
    \item \ADP{}-lite memory.
    \item Full \ADP{} memory.
\end{enumerate}

Then optionally compare retrieval-control variants:
\begin{enumerate}[label=\textbf{C\arabic*:},leftmargin=*,itemsep=2pt,topsep=2pt]
    \item Always retrieve.
    \item Fixed interval retrieve.
    \item Heuristic retrieve.
    \item Hierarchical trigger.
\end{enumerate}

\subsection{Metrics}
Define each metric carefully.

\paragraph{Task success.}
Binary or graded completion metric appropriate to the task.

\paragraph{Retrieval precision@k or candidate usefulness.}
Whether retrieved items are actually relevant for the live subproblem.

\paragraph{Injection usefulness.}
Whether the injected artifact measurably improves downstream behavior.

\paragraph{Injection harm rate.}
Fraction of injection events that degrade task progress, waste tokens, or mislead the agent.

\paragraph{Search cost and context cost.}
Latency, number of retrieval calls, tokens injected, and extra context budget consumed.

\paragraph{Robustness metrics.}
Performance under noisy, stale, contradictory, or weakly matching memories.

\subsection{Statistical testing}
State confidence intervals, number of seeds / runs, and statistical tests.
Examples:
paired bootstrap, paired permutation test, stratified confidence intervals by task family.

% =========================================================
% 7. Main Results
% =========================================================
\section{Main Results}

\subsection{Does \ADP{} improve retrieval-augmented performance?}
Present the main results table here.

% Example skeleton:
\begin{table}[t]
    \centering
    \small
    \caption{Main comparison across memory representations. Higher is better except harm and cost.}
    \label{tab:main-results}
    \begin{tabular}{lccccc}
        \toprule
        Method & Success $\uparrow$ & Usefulness $\uparrow$ & Harm $\downarrow$ & Searches $\downarrow$ & Cost $\downarrow$ \\
        \midrule
        No memory & -- & -- & -- & -- & -- \\
        Flat-text & -- & -- & -- & -- & -- \\
        Message-only & -- & -- & -- & -- & -- \\
        Summary & -- & -- & -- & -- & -- \\
        \ADP{}-lite & -- & -- & -- & -- & -- \\
        \ADP{} & -- & -- & -- & -- & -- \\
        \bottomrule
    \end{tabular}
\end{table}

\paragraph{What to write here.}
Answer RQ1 directly in the first paragraph.
Do not merely describe the table.
Say whether \ADP{} wins, by how much, and on which task families.

\subsection{Where does \ADP{} help most?}
Break results down by:
\begin{itemize}[leftmargin=*,itemsep=2pt,topsep=2pt]
    \item coding vs browsing vs mixed tasks,
    \item short vs long-horizon tasks,
    \item repeated subproblems vs novel tasks.
\end{itemize}

\subsection{When does retrieval hurt?}
This subsection is important for NeurIPS E\&D positioning.
Show:
\begin{itemize}[leftmargin=*,itemsep=2pt,topsep=2pt]
    \item cases where all retrieval hurts,
    \item cases where \ADP{} hurts less,
    \item cases where structure helps only when retrieval is selective.
\end{itemize}

% =========================================================
% 8. Ablations and Analysis
% =========================================================
\section{Ablations and Analysis}

\subsection{Which parts of \ADP{} matter?}
Ablate:
\begin{itemize}[leftmargin=*,itemsep=2pt,topsep=2pt]
    \item remove action types,
    \item remove tool kwargs,
    \item remove observations,
    \item remove metadata,
    \item reorder or collapse steps,
    \item control for token count.
\end{itemize}

\begin{table}[t]
    \centering
    \small
    \caption{Ablation of structured components within \ADP{} memory.}
    \label{tab:ablations}
    \begin{tabular}{lcccc}
        \toprule
        Variant & Success $\uparrow$ & Usefulness $\uparrow$ & Harm $\downarrow$ & Cost $\downarrow$ \\
        \midrule
        Full \ADP{} & -- & -- & -- & -- \\
        w/o action types & -- & -- & -- & -- \\
        w/o observations & -- & -- & -- & -- \\
        w/o tool kwargs & -- & -- & -- & -- \\
        w/o metadata & -- & -- & -- & -- \\
        token-matched flat text & -- & -- & -- & -- \\
        \bottomrule
    \end{tabular}
\end{table}

\subsection{Representation vs retrieval-control}
If you include control policies, this is where they belong.
Make the interaction explicit:
does \ADP{} still help under always-retrieve?
does its gain increase under selective search?

\subsection{Noise and staleness stress tests}
Corrupt the memory bank and measure robustness:
\begin{itemize}[leftmargin=*,itemsep=2pt,topsep=2pt]
    \item irrelevant memories,
    \item stale trajectories,
    \item conflicting procedures,
    \item partial trajectories.
\end{itemize}

\subsection{Qualitative case studies}
Provide 2--4 short examples:
\begin{enumerate}[leftmargin=*,itemsep=2pt,topsep=2pt]
    \item one success due to preserved procedure,
    \item one failure due to misleading lexical similarity,
    \item one failure of \ADP{} itself,
    \item one case where no memory was best.
\end{enumerate}

% =========================================================
% 9. Optional Formal Retrieval-Control Section
% =========================================================
\section{Optional: Hierarchical Retrieval-Control Formulation}
% Keep this section concise unless experiments directly support it.
% If the paper becomes too crowded, move most of this to the appendix.

We optionally formalize retrieval-control as a hierarchical decision process.
At step $t$, the wrapper observes a local state $x_t$ and derives a compact query fingerprint $F_t$.
A first-stage policy decides whether to query memory.
If retrieval is triggered, the system obtains candidate memories
\[
C_t = \{c_1,\dots,c_k\},
\]
scores them, and then chooses one of the actions
\[
\mathcal{A} = \{\texttt{NoSearch},\ \texttt{SearchNoInject},\ \texttt{InjectSummary},\ \texttt{InjectFull}\}.
\]

The reason for this decomposition is that some relevance quantities are only observable after retrieval.
Thus, the first-stage policy is not a full utility estimator; it is only a search gate.

\paragraph{Recommendation.}
If this section does not have strong experimental support, keep only one paragraph in the main paper and
move the equations to the appendix.

% =========================================================
% 10. Discussion
% =========================================================
\section{Discussion}

\subsection{What the results mean}
Interpret the results at the representation level, not just the system level.
Examples:
\begin{itemize}[leftmargin=*,itemsep=2pt,topsep=2pt]
    \item If \ADP{} wins mostly on procedural tasks, say that its value lies in preserving operational structure.
    \item If summary memory is competitive on easy tasks, say that structure is not always necessary.
    \item If \ADP{} helps retrieval but not injection, separate the two bottlenecks explicitly.
\end{itemize}

\subsection{Implications for future memory systems}
Argue that trajectory representation is a first-class design choice in agent memory.
Future work should not only optimize retrievers and rerankers, but also ask what kind of memory object is
worth retrieving in the first place.

% =========================================================
% 11. Limitations
% =========================================================
\section{Limitations}
This paper has several limitations.
First, the current study may focus on one primary agent framework and one specific memory backend, which
limits generalization across frameworks and storage systems.
Second, the present \Recall{} implementation may represent each session as a single indexed item, which is
coarser than segment-level or subtrajectory memory.
Third, the evaluation may not fully disentangle representation quality from query construction quality.
Fourth, any trigger or injection mechanism studied here should be understood as an initial policy rather
than a definitive optimal controller.
Finally, the tasks considered in this paper may emphasize procedural memory more than all real-world agent
deployments do.

% =========================================================
% 12. Reproducibility Statement
% =========================================================
\section{Reproducibility Statement}
We will release the code, prompts, configuration files, conversion scripts, trajectory schemas,
evaluation harness, and analysis scripts needed to reproduce the reported results.
The released artifact will include:
(i) conversion from raw session logs to \ADP{},
(ii) memory construction for all compared representations,
(iii) retrieval and evaluation scripts,
(iv) exact task splits,
(v) metric implementations, and
(vi) instructions for reproducing the main tables and figures.

% =========================================================
% 13. Conclusion
% =========================================================
\section{Conclusion}
This paper asked whether \ADP{}, originally proposed as a standardized format for offline agent
training data, also serves as a superior memory representation for retrieval-augmented agents.
By treating trajectory representation as an evaluative variable rather than a fixed implementation detail,
we showed how structured procedural memory can be studied rigorously.
Our results provide evidence about when structure helps, which components matter, and when retrieval is
beneficial versus harmful.
More broadly, the paper argues that representation choice is a central part of the science of agent
memory, not merely an engineering convenience.

% =========================================================
% References
% =========================================================
\bibliographystyle{plainnat}
\bibliography{references}

% =========================================================
% Appendix
% =========================================================
\appendix

\section{Appendix Overview}
This appendix provides implementation details, extended experiments, additional qualitative examples,
and optional mathematical material omitted from the main paper for space.

\section{Implementation Details}
\subsection{Current \Recall{} repository}
List only the modules you truly have.
For example:
\begin{itemize}[leftmargin=*,itemsep=2pt,topsep=2pt]
    \item configuration,
    \item session-to-\ADP{} conversion,
    \item embedding,
    \item vector-store ingestion,
    \item offline ingestion scripts.
\end{itemize}

\subsection{Online evaluation harness}
Describe any code added beyond the current repository for the experiments.

\section{Prompting, Models, and Runtime Configuration}
Specify:
\begin{itemize}[leftmargin=*,itemsep=2pt,topsep=2pt]
    \item base model(s),
    \item decoding settings,
    \item tool environment,
    \item retrieval top-$k$,
    \item summary budgets,
    \item context-window budgets,
    \item seed management.
\end{itemize}

\section{Metric Definitions}
Give exact formulas or pseudocode for:
\begin{itemize}[leftmargin=*,itemsep=2pt,topsep=2pt]
    \item retrieval usefulness,
    \item injection harm,
    \item search cost,
    \item robustness metrics.
\end{itemize}

\section{Extended Tables}
Put full per-task-family tables here.

\section{Additional Case Studies}
Show longer examples of retrieved trajectories and injection outcomes.

\section{Optional Formalization of the Trigger}
If retained, move the full trigger equations here.

\subsection{Local state and fingerprint}
\[
x_t = (g_t, h_t, a_t, e_t, d_t, r_t), \qquad F_t = \phi(x_t).
\]

\subsection{Candidate retrieval}
\[
C_t = \{c_1,\ldots,c_k\}.
\]

\subsection{Candidate scoring}
For each candidate $c_i$, define:
\[
s_i \in [0,1],\quad
v_i \in [0,1],\quad
m_i \in [0,1],\quad
q_i \in [0,1],\quad
p_i \in [0,1].
\]
A generic ranking rule may be written as
\[
R_i = \alpha s_i + \beta v_i + \gamma m_i + \delta q_i - \lambda p_i,
\qquad
R_t^\star = \max_i R_i.
\]
Keep this generic unless you have experimental justification for fixed weights.

\subsection{Search gate}
If you keep the local gate:
\[
L_t = \Gamma_d(t)\bigl(w_E E_t + w_N N_t^{\text{loc}} + w_U U_t\bigr),
\]
with a hard-trigger override and a feasibility constraint for budget.
Again, avoid overclaiming any fixed constants unless validated.

\section{Artifact Checklist for Release}
\begin{itemize}[leftmargin=*,itemsep=2pt,topsep=2pt]
    \item code repository,
    \item dependency file,
    \item scripts to regenerate memory banks,
    \item evaluation runner,
    \item seed list,
    \item anonymized sample data or instructions for reconstruction,
    \item README with exact commands,
    \item license statement.
\end{itemize}

% =========================================================
% End of document
% =========================================================
\end{document}